% Part: propositional-logic
% Chapter: propositional-logic
% Section: preliminaries

\documentclass[../../../include/open-logic-section]{subfiles}

\begin{document}

\olfileid[zh]{pl}{syn}{pre}

\olsection{预备知识}

\begin{thm}[\emph{关于!!{formula}的归纳原理}]
  \ollabel{thm:induction}
  若某个性质~$P$ 对所有原子!!{formula}成立，且满足：
  \begin{enumerate}
    \tagitem{prvNot}{每当它对~$!A$ 成立，它就对 $\lnot !A$ 成立；}{}
    \tagitem{prvAnd}{每当它对 $!A$ 与~$!B$ 成立，它就对 $(!A \land !B)$ 成立；}{}
    \tagitem{prvOr}{每当它对 $!A$ 与~$!B$ 成立，它就对 $(!A \lor !B)$ 成立；}{}
    \tagitem{prvIf}{每当它对 $!A$ 与~$!B$ 成立，它就对 $(!A \lif !B)$ 成立；}{}
    \tagitem{prvIff}{每当它对 $!A$ 与~$!B$ 成立，它就对 $(!A \liff !B)$ 成立；}{}
  \end{enumerate}
  则 $P$ 对所有!!{formula}成立。
\end{thm}

\begin{proof}
  设 $S$ 为所有具有性质~$P$ 的!!{formula}之全体。显然 $S \subseteq \Frm[L_0]$。$S$~满足 \olref[fml]{defn:formulas} 的全部条件：它包含所有原子!!{formula}，并且对所有!!{operator}封闭。$\Frm[L_0]$ 是满足这些条件的最小类，因此 $\Frm[L_0] \subseteq S$。于是 $\Frm[L_0] = S$，即每个公式都具有性质~$P$。
\end{proof}

\begin{prop}\ollabel{prop:balanced}
   在~$\Frm[L_0]$ 中任意一个!!{formula}是\emph{平衡}的，即其左括号与右括号数目相同。
\end{prop}

\begin{prob}
证明 \olref[pl][syn][pre]{prop:balanced}
\end{prob}

\begin{prop} \ollabel{prop:noinit}
任意!!a{formula}的真初始段都不是!!a{formula}。
\end{prop}

\begin{prob}
  证明 \olref[pl][syn][pre]{prop:noinit}
\end{prob}

\begin{prop}[唯一可读性]
在 $ \Frm[L_0]$ 中任意一个!!{formula}~$!A$ 恰有如下一种解析：
\begin{enumerate}
\tagitem{prvFalse}{$\lfalse$。}{}

\tagitem{prvTrue}{$\ltrue$。}{}

\item 对于某个 $\Obj p_n \in  \PVar$，$\Obj p_n$。

\tagitem{prvNot}{对于某个!!{formula}~$!B$，$\lnot !B$。}{}

\tagitem{prvAnd}{对于某些!!{formula} $!B$ 与~$!C$，$(!B \land !C)$。}{}

\tagitem{prvOr}{对于某些!!{formula} $!B$ 与~$!C$，$(!B \lor !C)$。}{}

\tagitem{prvIf}{对于某些!!{formula} $!B$ 与~$!C$，$(!B \lif !C)$。}{}

\tagitem{prvIff}{对于某些!!{formula} $!B$ 与~$!C$，$(!B \liff !C)$。}{}
\end{enumerate}
而且，这种解析是\emph{唯一}的。
\end{prop}

\begin{proof}
对 $!A$ 归纳。例如，设 $!A$ 分别以 $(!B \lif !C)$ 与 $(!B' \lif !C')$ 两种不同方式解析。则 $!B$ 与 $!B'$ 必相同（否则其中之一是另一个的真初始段）；因此若 $!A$ 的两种解析不同，则必是因为 $!C$ 与 $!C'$ 是同一种符号序列的两种不同的 $\lif$ 后件解析，而由归纳假设这是不可能的。
\end{proof}


\begin{defn}[统一代入]
若 $!A$ 与 $!B$ 是!!{formula}，且 $\Obj p_i$ 是一个 !!{propositional
variable}，则 $\Subst{!A}{!B}{\Obj p_i}$ 表示在 $!A$ 中把 $\Obj p_i$ 的每一处出现替换为 $!B$ 的一处出现所得的结果；类似地，把 $\Obj p_1$、\dots、~$\Obj p_n$ 同时替换为!!{formula} $!B_1$、\dots、~$!B_n$ 记为
$\SSubst{!A}{\subst{!B_1}{\Obj p_1},\dots,\subst{!B_n}{\Obj p_n}}$。
\end{defn}

\begin{prob} 对下面五个!!{formula}中的每一个，判断该!!{formula}是否能表示为代入 \( \Subst{!A}{!B}{\Obj p_i} \),
 其中 \( !A \) 分别是 (i) \( \Obj p_0 \)；(ii) \( ( \lnot \Obj p_0 \land \Obj
 p_1) \)；(iii) \( ( ( \lnot \Obj p_0 \lif \Obj p_1 ) \land \Obj
 p_2 ) \)。在每种情形下指明相应的代入。
  \begin{enumerate}
    \item \( \Obj p_1 \)
    \item \( ( \lnot \Obj p_0 \land \Obj p_0 ) \)
    \item \( ( ( \Obj p_0 \lor \Obj p_1 ) \land \Obj p_2 )  \)
    \item \( \lnot ( ( \Obj p_0 \lif \Obj p_1 ) \land \Obj p_2 ) \)
    \item \( (( \lnot ( \Obj p_0 \lif \Obj p_1 ) \lif ( \Obj p_0 \lor \Obj p_1 )) \land \lnot ( \Obj p_0 \land \Obj p_1 )) \)
  \end{enumerate}
\end{prob}

\begin{prob}
用归纳给出 $\Subst{!A}{!B}{p}$ 的一个数学上严格的定义。
\end{prob}

\end{document}
